\begin{thesisabstract}

Reconstructed 3D scenes are, almost without exception, empty. A gym is captured without anyone running on its treadmill, a kitchen without anyone opening its appliances: the rooms are faithful, but the people who give them purpose are missing. Placing a person into such a scene, caught in the act of using it, demands more than a body that looks right. The interaction must be semantically correct, performing the requested action on the right object; physically valid, balanced and free of intersection; and in true contact with the scene, each hand meeting the handle it grips, each foot landing on the surface that supports it. These demands are not equally hard. Pretrained foundation models have absorbed what human activity looks like and supply the semantics readily, and physical validity can be enforced directly once body and scene are both expressed as geometry. Contact is what neither provides, and it is what ties an interaction to the specific scene it is placed in rather than to a generic image of the activity: which region of which surface, precisely, each participating body part must touch. Prior work leaves this question open. Methods that learn contact from captured data do not transfer beyond their training scenes; language-driven control decides which object is engaged, but not where on its surface contact falls. The closest zero-shot methods resolve the scene side only partially: some infer contact implicitly from a generated image, an inference that may carry a simple support interaction but strains when several body parts must each find their own precise hold, as when a person steps down from a ladder, and a contact that is wrong or missing can then be neither detected nor corrected; others commit to a named scene element in advance, leaving the contact location on that element undetermined.

This thesis presents a zero-shot method that synthesizes a static human--scene interaction from a natural-language instruction, a single posed RGB view, and the reconstructed scene mesh, producing an SMPL-X body in the scene's metric coordinate frame using only pretrained foundation models, explicit camera geometry, and optimization, without any training on paired interaction data. The key idea is to use foundation models to localize where the interaction touches the scene: generation supplies the interaction knowledge, and camera geometry and optimization turn it into a physically correct body. The central contribution is an agentic contact-estimation procedure: an image-generation model proposes body-part-specific contact regions, which are composited onto the unmodified posed view to preserve exact correspondence with the calibrated camera, and a vision-language evaluator inspects each composite against the intended interaction, issuing corrections in a closed loop until the contacts are accepted. The verified regions are projected onto the mesh as explicit 3D contact targets, and a scene-aware optimization refines a monocularly initialized body to meet them while avoiding penetration.

The method is evaluated on 23 interactions across 15 reconstructed ScanNet++ scenes, against manually annotated scene-side contact references and a vision-language rubric introduced for this setting. The verify-and-correct loop eliminates missed contacts entirely and nearly halves the final contact error relative to accepting a single generation, while leaving penetration unchanged, so its effect falls precisely where the loop is designed to act: on the accuracy of the scene-side contact. Compared with PhySIC, the closest single-image human--scene contact baseline, the method places the contacting body parts markedly more precisely on the intended surface regions while penetrating the scene less. These results demonstrate that recovering precise, verified, per-body-part contact is what turns the image-space interaction knowledge of foundation models into physically grounded 3D humans, without task-specific training.

\end{thesisabstract}